\scanpage{55}
\mspage{14}
\uline{Dém} \deleted{Tout ceci}

$E$ est engendré par $R$, et même par $r_1\ r_2\ r_3\ r_4\ r_5\ r_6$

Car des sommes partielles des $r_i$ \inserted{$1\le i\le5$} donnent les $\zeta_{i}-\zeta_j$ $(i<j)$, et il est clair que les $\zeta_i-\zeta_1$ $(1\le i\le6)$ et \deleted{$\unsure{\eta-\lambda-2\zeta_1|\eta-2\lambda}=\eta-\lambda$ (car $\lambda-\eta=r_0=r'_6$} \inserted{$\eta-3\zeta_1$} engendrent $E$ $(=\operatorname{Ker}(\widehat E\xrightarrow{\lambda\, .}\mathbb Z))$ puisque \deleted{$\lambda.\zeta_1=1$, et} $\lambda.\zeta_i=1,\ \lambda.\eta=3$. Donc $E$ engendré
\deleted{$\eta,\zeta_1,\ldots,\zeta_6$ engendrent $\widehat E$ et leurs produits avec $\lambda$ sont $3\ 1$}
par \inserted{les} $\zeta_1-\zeta_i$ et \deleted{$\eta-3\zeta_1=3\zeta_1-\eta$} \inserted{$\eta-3\zeta_1$}
\deleted{\unsure{$= \lambda+\eta$|$3\lambda-\eta$}} donc par $r_1\ldots r_5$ et $\eta-3\zeta_1$. Reste à prouver que $\eta-3\zeta_1$ est une comb. des $r_i$, \unsure{rappeler|rappelons}
\[
\begin{aligned}
r_6&=\deleted{\unsure{r_0=2\eta}}\ \eta-\zeta_1-\zeta_2-\zeta_3,\\
&=(\eta-3\zeta_1)+(\zeta_1-\zeta_2)+
\underset{\substack{\Vert\\(\zeta_1-\zeta_2)+(\zeta_2-\zeta_3)}}{(\zeta_1-\zeta_3)},\quad\text{\unsure{gagné|on gagne}.}
\end{aligned}
\]

On sait déjà que les $s_r$ transforment $\Delta$ en lui $=$, donc $R$ en lui-$=$, et comme $s_r$\deleted{\unsure{$=(r\mapsto-r)$}} \inserted{$=(x\longmapsto x+r^\star(x)r)$ avec $r^\star(r)=-2$,} on conclut que $R$ est un syst. de racines --- compte tenu que deux racines distinctes \unsure{non opp.} sont \inserted{non} proportionelles. \deleted{Que} \inserted{Prouvons que} $\alpha_1\ldots\alpha_6$ sont un syst. de racines simples. \deleted{\unsure{on voit en}}

On a déjà vu
\[
\boxed{r_{ij}=r_i+r_{i+1}\ldots+r_{j-1}\qquad\text{si }i<j}
\quad\text{(\unsure{pr $i$ ou|pris en|pour $1\le i,$} }j=\unsure{5|i},\quad\text{\unsure{max.|max|$\le$} : }5\text{)}
\]
On a
\[
\begin{aligned}
r_{ijh}&=\eta-\zeta_i-\zeta_j-\zeta_h=(\eta-\zeta_1-\zeta_2-\zeta_3)
+\underbrace{(\zeta_1-\zeta_i)}_{r_{1i}}+\underbrace{(\zeta_2-\zeta_j)}_{r_{2j}}+\underbrace{(\zeta_3-\zeta_h)}_{r_{3h}}\\
&=\deleted{\unsure{\Sigma_6+}}\ r_{123}+r_{1i}+r_{2j}+r_{\unsure{2|3}h}
\end{aligned}
\]
donc
\[
r_{ijh}=r_6+(r_1\ldots+r_{i-1})+(r_2\ldots+r_{j-1})+(r_3+\ldots r_{h-1})
\]
% archive sheet number: 54
